\documentclass[10pt,a4paper]{article}
\usepackage[a4paper,margin=14mm,top=12mm,bottom=14mm]{geometry}
\usepackage{fontspec}
\setmainfont{Liberation Sans}
\usepackage{polyglossia}
\setdefaultlanguage{french}
\usepackage{amsmath}
\usepackage{xcolor}
\usepackage{graphicx}
\usepackage[most]{tcolorbox}
\usepackage{tabularx,booktabs,array,multirow}
\usepackage{enumitem}
\usepackage{fancyhdr}
\usepackage{lastpage}
\usepackage[hidelinks]{hyperref}
\definecolor{uimm}{HTML}{005677}
\definecolor{cyan}{HTML}{13A2DD}
\definecolor{fond}{HTML}{E6F2FA}
\definecolor{ok}{HTML}{2E7D32}
\definecolor{ko}{HTML}{B3261E}
\definecolor{violet}{HTML}{7B1FA2}
\setlength{\parindent}{0pt}
\setlength{\parskip}{3pt}
\setlist{nosep}
\setlist[itemize]{leftmargin=13pt,label=\textcolor{cyan}{\textbullet}}
\newcommand{\titre}[1]{\vspace{5pt}{\color{uimm}\large\bfseries #1}\par\vspace{-3pt}{\color{cyan}\rule{\linewidth}{1.2pt}}\par\vspace{0pt}}
\tcbset{boite/.style={colback=fond,colframe=uimm,boxrule=0.6pt,arc=2pt,left=6pt,right=6pt,top=4pt,bottom=4pt},
        defi/.style={colback=violet!6,colframe=violet,boxrule=0.8pt,arc=2pt,left=6pt,right=6pt,top=4pt,bottom=4pt,
                     title={\bfseries Défi},coltitle=white,colbacktitle=violet,fonttitle=\small}}
\newcommand{\bandeau}[3]{%
\begin{tcolorbox}[colback=uimm,colframe=uimm,arc=2pt,left=4pt,right=8pt,top=4pt,bottom=4pt]
\begin{minipage}[c]{18mm}\setlength{\fboxsep}{1.2pt}\fcolorbox{white}{white}{\includegraphics[height=16.5mm]{logo_uimm.jpg}}\end{minipage}\hspace{4mm}%
\begin{minipage}[c]{\dimexpr\linewidth-22mm\relax}\color{white}
{\small #1}\\[1pt]{\LARGE\bfseries #2}\\[2pt]{\large #3}\hfill{\small CQPM Technicien qualité}
\end{minipage}
\end{tcolorbox}}
\newcommand{\reponse}[1][1]{\par\vspace{2pt}{\color{gray!60}\foreach \i in {1,...,#1}{\rule{\linewidth}{0.3pt}\par\vspace{9pt}}}}
\usepackage{pgffor}
\pagestyle{fancy}\fancyhf{}
\renewcommand{\headrulewidth}{0pt}
\newcommand{\pied}[1]{\fancyfoot[L]{\footnotesize\color{uimm}#1}\fancyfoot[R]{\footnotesize\color{uimm}Page \thepage{} sur \pageref{LastPage}}}

\usepackage{tikz}
\usepackage{longtable}
\renewcommand{\bandeau}[3]{%
\begin{tcolorbox}[colback=uimm,colframe=uimm,arc=2pt,left=4pt,right=8pt,top=4pt,bottom=4pt]
\begin{minipage}[c]{18mm}\setlength{\fboxsep}{1.2pt}\fcolorbox{white}{white}{\includegraphics[height=16.5mm]{logo_uimm.jpg}}\end{minipage}\hspace{4mm}%
\begin{minipage}[c]{\dimexpr\linewidth-22mm\relax}\color{white}
{\small #1}\\[1pt]{\LARGE\bfseries #2}\\[2pt]{\large #3}
\end{minipage}
\end{tcolorbox}}
% en-tête d'une fiche : code, titre, question, bloc, durée, organisation, niveau
\newcommand{\entete}[7]{%
\begin{tcolorbox}[colback=uimm,colframe=uimm,arc=2pt,left=4pt,right=6pt,top=4pt,bottom=4pt]
\begin{minipage}[c]{18mm}\setlength{\fboxsep}{1.2pt}\fcolorbox{white}{white}{\includegraphics[height=16.5mm]{logo_uimm.jpg}}\end{minipage}\hspace{3mm}%
\begin{minipage}[c]{\dimexpr\linewidth-58mm\relax}\color{white}\raggedright
{\small Base SPC · #4}\\[1pt]{\Large\bfseries #1 · #2}\\[2pt]{\itshape #3}
\end{minipage}\hfill
\begin{minipage}[c]{33mm}\color{white}\small\raggedleft
Durée : \textbf{#5}\\ \textbf{#6}\\ Niveau #7 sur 3
\end{minipage}
\end{tcolorbox}}
\newcommand{\tagE}{{\setlength{\fboxsep}{1.5pt}\colorbox{cyan}{\color{white}\scriptsize\bfseries Ellistat}}}
\newcommand{\tagW}{{\setlength{\fboxsep}{1.5pt}\colorbox{uimm}{\color{white}\scriptsize\bfseries Page web}}}
\tcbset{pred/.style={colback=orange!6,colframe=orange!80!black,boxrule=0.6pt,arc=2pt,left=6pt,right=6pt,top=3pt,bottom=1pt,
                    title={\bfseries Je prédis, avant de commencer},coltitle=white,colbacktitle=orange!80!black,fonttitle=\small},
        retenir/.style={colback=fond,colframe=uimm,boxrule=0.6pt,arc=2pt,left=6pt,right=6pt,top=3pt,bottom=3pt,
                    title={\bfseries À retenir},coltitle=white,colbacktitle=uimm,fonttitle=\small},
        plus/.style={colback=violet!6,colframe=violet,boxrule=0.6pt,arc=2pt,left=6pt,right=6pt,top=3pt,bottom=3pt,
                    title={\bfseries Pour aller plus loin},coltitle=white,colbacktitle=violet,fonttitle=\small},
        act/.style={colback=uimm!10,colframe=uimm,boxrule=0.5pt,arc=2pt,left=5pt,right=5pt,top=2pt,bottom=2pt,before skip=7pt,after skip=3pt},
        verif/.style={colback=orange!6,colframe=orange!80!black,boxrule=0.6pt,arc=2pt,left=6pt,right=6pt,top=3pt,bottom=3pt,
                    title={\bfseries À vérifier sur poste avant la première séance},coltitle=white,colbacktitle=orange!80!black,fonttitle=\small}}
\newcommand{\fichier}[1]{\texttt{\small #1}}
\newcommand{\menu}[1]{\textsf{\textbf{#1}}}
\newcommand{\champ}[1]{\rule{#1}{0.3pt}}

\pied{Base SPC · Mémo · Les cartes aux attributs}
\begin{document}
\thispagestyle{fancy}
\bandeau{Base SPC · Mémo}{Les cartes aux attributs}{Je compte au lieu de mesurer : quelle carte choisir, et pourquoi ?}

\titre{1. L'idée en trois phrases}
Une carte aux attributs suit un \textbf{comptage} au fil du temps. Elle sert à voir si le procédé est stable ou s'il dérive. Pour choisir la carte, deux questions suffisent.

\begin{tcolorbox}[boite]
\textbf{Question 1. Qu'est-ce que je compte ?}\\
Des \textbf{pièces défectueuses} : la pièce est refusée, qu'elle porte un défaut ou cinq. Ou des \textbf{défauts} : chaque anomalie compte, plusieurs peuvent se trouver sur la même pièce.\par\smallskip
\textbf{Question 2. L'unité contrôlée a-t-elle toujours la même taille ?}\\
Même nombre de pièces par lot, même longueur de rouleau, même surface de capot : taille \textbf{fixe}. Sinon : taille \textbf{variable}.
\end{tcolorbox}

\begin{center}\renewcommand{\arraystretch}{1.6}
\begin{tabularx}{0.92\linewidth}{|>{\bfseries}p{38mm}|>{\centering\arraybackslash}X|>{\centering\arraybackslash}X|}\hline
 & \textbf{Taille fixe} & \textbf{Taille variable}\\\hline
Pièces défectueuses & \textbf{Np} (ou p) & \textbf{p}\\\hline
Défauts & \textbf{c} & \textbf{u}\\\hline
\end{tabularx}\end{center}

\begin{tcolorbox}[retenir]\textbf{Fixe : on compte. Variable : on divise.} Avec une taille fixe, le nombre brut se compare directement d'une fois sur l'autre (Np, c). Avec une taille variable, on le rapporte à la taille pour comparer ce qui est comparable (p, u).\end{tcolorbox}

\titre{2. Deux images pour s'en souvenir}
\begin{itemize}
\item \textbf{Pièce défectueuse, c'est un candidat recalé.} Qu'il ait raté une question ou dix, il compte pour un seul recalé. On compte des candidats : carte p ou Np.
\item \textbf{Défaut, c'est une faute sur une copie.} Une copie peut en porter plusieurs et chacune compte. On compte des fautes : carte c ou u.
\item \textbf{Taille variable : on compare des copies de 2 pages et de 10 pages.} Dix fautes sur dix pages ne valent pas dix fautes sur deux pages. On rapporte à la page : c'est la carte u.
\end{itemize}

\titre{3. Les calculs, avec un exemple de la base pour chacun}
\renewcommand{\arraystretch}{1.55}
\begin{tabularx}{\linewidth}{|>{\bfseries}c|>{\raggedright\arraybackslash}p{50mm}|>{\raggedright\arraybackslash}X|}\hline
\textbf{Carte} & \textbf{Ligne centrale et limites} & \textbf{Exemple chiffré}\\\hline
p & $\bar p = \sum \text{défectueuses} \,/\, \sum \text{pièces}$\par $\bar p \pm 3\sqrt{\bar p\,(1-\bar p)\,/\,n_i}$ & 35 lots de 50 pièces, 103 défectueuses : $\bar p = 5{,}89$ \%, limite supérieure 15,87 \%. Si $n_i$ change, la limite suit : ligne brisée.\\\hline
Np & $n\bar p \pm 3\sqrt{n\bar p\,(1-\bar p)}$\par (lots de taille $n$ fixe) & Mêmes lots : $n\bar p = 2{,}94$ pièce, limite supérieure 7,9 pièces (soit $50 \times 15{,}87$ \%).\\\hline
c & $\bar c \pm 3\sqrt{\bar c}$\par (unités identiques) & 25 panneaux, 111 défauts : $\bar c = 4{,}44$, limite supérieure $4{,}44 + 3\sqrt{4{,}44} = 10{,}76$.\\\hline
u & $\bar u = \sum \text{défauts} \,/\, \sum \text{tailles}$\par $\bar u \pm 3\sqrt{\bar u \,/\, n_i}$ & 25 rouleaux de 40 à 120 m : $\bar u = 0{,}055$ défaut par mètre. Rouleau 14 : 13 défauts sur 80 m, $u = 0{,}1625$, limite 0,134 : hors limites.\\\hline
\end{tabularx}
\par\smallskip
Si la limite inférieure calculée est négative, on la remplace par 0. Un point hors limites ou une série de points du même côté de la moyenne signale que \textbf{quelque chose a changé} : on le cherche avant de corriger quoi que ce soit.

\titre{4. Les cartes de Laney (P Laney et U Laney)}
Les cartes p et u supposent que les lots varient seulement par hasard, comme un tirage. Quand les lots sont très grands, les vrais écarts d'un lot à l'autre (matière, réglage, équipe) dépassent ce hasard. La carte classique montre alors des points rouges presque partout, et on ne sait plus lesquels comptent.\par
La correction de Laney multiplie la largeur des limites par un facteur $\sigma_z$, estimé à partir de l'écart d'un lot au suivant. \textbf{Si $\sigma_z$ est voisin de 1, la carte classique suffit.} S'il est nettement plus grand, les limites s'élargissent et seuls les points vraiment anormaux sortent. P Laney remplace la carte p ; U Laney remplace la carte u.

\titre{5. Dans Ellistat}
\begin{enumerate}[label=\textcolor{cyan}{\textbf{\arabic*.}},leftmargin=15pt,itemsep=1.5pt]
\item Menu \menu{Stats Desc.}, onglet \menu{Cartes de contrôle aux attributs}.
\item « Attributs » : la colonne des défectueuses ou des défauts. « Taille de lot » : la colonne des tailles, \textbf{si elles varient}.
\item Pour la carte c, taper \textbf{1} dans « Taille de lot (si fixe) » : sans cette valeur, Ellistat n'a tracé que 25 points sur 30.
\item « Phase I » : Début et Fin. \textbf{Vérifier la valeur de Fin}, qui garde celle de l'étude précédente. Ces lots fixent la ligne centrale et les limites.
\item Cliquer sur le bouton rond de la carte : Np, c, p, u, P Laney ou U Laney.
\end{enumerate}

\titre{6. Je m'entraîne}
\begin{enumerate}[label=\textcolor{cyan}{\textbf{Q\arabic*.}},leftmargin=22pt,itemsep=2pt]
\item 200 vis sont contrôlées chaque heure au calibre ; on note les vis refusées. Quelle carte ? \champ{20mm}
\item On compte les grains de peinture sur chaque capot, tous identiques. Quelle carte ? \champ{20mm}
\item Rouleaux de 40 à 120 m ; on note les défauts de chaque rouleau. Quelle carte ? \champ{20mm}
\item Lots reçus de 80 à 400 pièces ; on note les pièces refusées. Quelle carte ? \champ{20mm}
\item Une pièce porte trois rayures : elle compte pour \champ{8mm} sur une carte p et pour \champ{8mm} sur une carte c.
\end{enumerate}
\begin{tcolorbox}[colback=gray!8,colframe=gray!60,boxrule=0.5pt,arc=2pt,left=6pt,right=6pt,top=3pt,bottom=3pt]\footnotesize\textbf{Réponses.} Q1 Np (ou p) : défectueuses, taille fixe. Q2 c : défauts, unité constante. Q3 u : défauts, taille variable. Q4 p : défectueuses, taille variable. Q5 1 sur la carte p ; 3 sur la carte c.\end{tcolorbox}

\vfill{\footnotesize\color{uimm}\textbf{Pour continuer dans la base :} fiches A1 (choisir la carte), A2 à A4 (carte p), A5 (p ou Np), A6 (carte c), A7 (carte u), A8 (période de référence et P Laney).}
\end{document}
